\section{May 14}

Last time, we have introduced the higher Chow groups as the homology of the complex $Z^p(X, \bullet)$,
where $Z^p(X, n)$ is the free abelian group generated by the codimension $p$ subvarieties of $X \times \Delta^n$ that intersect all faces properly.

\subsection{}

  \begin{theorem}
    Let $X$ be a scheme (of finite type over a field) and $p_X: X \times \bbA^1 \to X$ be the projection. 
    Then $p_X^*: Z^i(X, \bullet) \to Z^i(X \times \bbA^1, \bullet)$ is a quasi-isomorphism, i.e., it induces an isomorphism on homology:
    \[ \CH^i(X, n) \cong \CH^i(X \times \bbA^1, n) \quad \text{for all } n. \]
  \end{theorem}
  \begin{proof}
    For $n \geq 0$, identify $\bbA^1 \times \Delta^n$ with $\Delta^{n+1}$ and define $T_n: Z^i(X \times \bbA^1, n) \to Z^i(X, n+1)$ s.t. 
    \[ \partial T_n + T_{n-1} \partial = i_0^* - i_1^*, \]
    where $i_0, i_1: X \to X \times \bbA^1$ are the inclusions at $0$ and $1$ respectively. 
    Then $T_\bullet$ gives a homotopy between $i_0^*$ and $i_1^*$.

    Now let $\tau: X \times \bbA^1 \times \bbA^1 \to X \times \bbA^1$ be the map defined by $\tau(x, t, s) = (x, st)$, which is flat.
    Define 
    \[ I_t: X \times \bbA^1 \injmap X \times \bbA^1 \times \bbA^1, \quad I_t(x, s) = (x, s, t). \]
    Then $\tau \circ I_0 = i_0$ and $\tau \circ I_1 = \id_{X \times \bbA^1}$.
    Since $I_t$ is $i_t$ for $X \times \bbA^1$, we have $I_0^*$ is homotopic to $I_1^*$.
    Hence $i_0^* = I_0^* \circ \tau^* \sim I_1^* \circ \tau^* = \id_{Z^i(X \times \bbA^1, \bullet)}$.

    \Yang{the argument is wrong, need to be fixed.}
    \Yang{no canonical identification between $\bbA^1 \times \Delta^n$ and $\Delta^{n+1}$,}
    \Yang{$T_n$, $i_0^*$, $i_1^*$ are not defined on $Z^i(X \times \bbA^1, n)$, but on a quasi-isomorphic subcomplex.}

    Recall the $\Delta^n = \Spec k[t_0, \dots, t_n]/(\sum t_i - 1)$.
    The vertices of $\bbA^1 \times \Delta^n$ are given by $(0, e_i)$ and $(1, e_i)$ for $i = 0, \dots, n$, where $e_i$ is the $i$-th vertex of $\Delta^n$.
    Identify $\bbA^1 \times \Delta^n$ with $\bbA^{n+1}$ by $(s, t_0, \dots, t_n) \mapsto (s, t_1, \dots, t_n)$.
    % Then the faces of $\bbA^1 \times \Delta^n$ are given by $s = 0$, $s = 1$, and $t_i = 0$ for $i = 0, \dots, n$.
    For a finite collection of vertices $S=\{v_0, \dots, v_r\}$ of $\bbA^1 \times \Delta^n$, define $\theta_S: \Delta^r \injmap \bbA^1 \times \Delta^n$ by $\theta_S(t_0, \dots, t_r) = \sum t_i v_i$.
    A triangulation of $\bbA^1 \times \Delta^n$ is a collection $\{\calT_n\}$, where each $\calT_n$ is a finite collection of finite set $S$ of vertices  of $\bbA^1 \times \Delta^n$ with $\sigma(S) \in \{\pm\}$ such that 
    \[ T_n := \sum_{S \in \calT_n} \sigma(S) \theta_S^*: Z^i(X \times \bbA^1, n) \to Z^i(X, n+1) \]
    satisfies $(\partial T_n + T_{n-1} \partial)(\alpha) =  \alpha|_{X\times \{0\} \times \Delta^n} - \alpha|_{X\times \{1\} \times \Delta^n}$ whenever defined.

    \Yang{}
  \end{proof}





